\cotefile{157-5}
\scanpage{309}
\mspage{225}
% archive sheet number: 308

2\textsuperscript{\unsure{o}}) Soit $\unsure{\pi|\kappa}$ un cardinal infini, et considérons les conditions :

$a_\pi)$ $I$ \unsure{grde|gde} devant $\pi$, $f$ cofinal, i.e. $\mathrm{F}1$, $\mathrm{F}2$.

$\deleted{\unsure{b}}c_\pi)$ $I'$ grde devant $\pi$, et \deleted{satisfait} \inserted{$f$ satisfait $\mathrm{F}1$ et}

\tikz[baseline=(f.base)]\node[draw,ellipse,inner sep=2pt](f){$\mathrm{F}3_\pi$}; Pour \deleted{\unsure{toute}} famille $(i_\alpha)_{\alpha\in A}$ d'objets de $I$, avec $\operatorname{card} A\leq\pi$, et $i\in\mathrm{Ob}\,I$ et d'une famille de flèches $u'_\alpha:f(\unsure{i'_\alpha|i_\alpha})\mathbin{\unsure{-|\longrightarrow}}f(i)$, existe $j\in\mathrm{Ob}\,I$ et une flèche $v':f(i)\mathbin{-}f(j)$, telles que $v'$ \uline{et} les $\unsure{v'_\alpha u'_\alpha|v'u'_\alpha}$ proviennent de flèches de $I$, i.e. \deleted{\unsure{il|ils}} existe des flèches $v:i\longrightarrow j$, $w_\alpha:i_\alpha\mathbin{-}j$ telles que

\[f(v)u'_\alpha=f(w_\alpha)\qquad\forall\,\alpha\in A.\]

\margin{\diagram{\begin{tikzcd}[ampersand replacement=\&,column sep=small,row sep=small]
i_\alpha\arrow[rr,dashed,bend left=18,"w_\alpha"]\arrow[d,no head]\& i\arrow[r,dotted,no head,"v"]\arrow[d,no head]\&j\arrow[d,no head]\\
f(i_\alpha)\arrow[r,"u'_\alpha"]\& f(i)\arrow[r,"v'=f(v)"]\&f(j)
\end{tikzcd}}}

% slanted
\margin{\uline{NB} La condition $\mathrm{F}3_\pi$ est \unsure{tautologiquement} satisfaite si $f$ est \uline{pleinement} \uline{fidèle}.}

$b_\pi)$ $I'$ grde devant $\pi$, les $ {}_{i'}\backslash I$ $(i'\in\mathrm{Ob}\,I')$

$b''_\pi)$ \inserted{\unsure{''\quad''|aussi|ainsi}} $I'$ grde devant $\pi$, les $ {}_{i'}\backslash I$ $(i'\in\mathrm{Ob}\,I')$ sont \inserted{\unsure{$\pi$-}} \uline{\unsure{pseudo-filtr.tes|pseudo-filtrantes}} \inserted{\unsure{i.e. satisfont} $\mathrm{PS}1_\pi'$}

\deleted{\unsure{On a}} \unsure{A|À} priori, on a

\diagram{\begin{tikzcd}[ampersand replacement=\&,row sep=tiny,column sep=small]
\& \deleted{\unsure{b_\pi}} \& \&\\
a_\pi\arrow[ur,Rightarrow,"\diagup"{pos=.7,inner sep=0pt}]\arrow[r,Leftrightarrow]\arrow[dr,Rightarrow,dash pattern=on 1pt off 1pt,"\diagup"{pos=.53,inner sep=0pt}]\& b_\pi\arrow[r,Rightarrow]\&c_\pi\arrow[r,Rightarrow]\&b''_\pi\\
\&\deleted{\unsure{c_\pi}}\&\&
\end{tikzcd}}

et si $f$ est fidèle, \inserted{\unsure{ou $I$ ordonnée},} les \deleted{trois} \inserted{quatre} conditions $a_\pi,b_\pi,\inserted{b''_\pi,}c_\pi$ sont équivalentes.

% slanted
\margin{\unsure{Dém de la proposition} \unsure{donnée en|dernière en} \unsure{cf $4.9.6.7$|$4.9.6.7$}, page 227.}

\uline{4.9.6.4.}

\uline{NB} Dans le cas qui nous intéresse, où $f$ est le foncteur \unsure{les|des|cas} (p.~223), on sait que $I'$ est grde devant $\pi$. On voudrait établir que $I$ l'est et que $f$ est cofinal, i.e. la condition $a_\pi)$. On \deleted{\unsure{applique}} le prouvera en prouvant $b_\pi$ (qui implique $a_\pi$, puisque $f$ est fidèle). Il reste donc à

\scanpage{310}
\mspage{226}
% archive sheet number: 309

prouver, sur le foncteur $f$ du cas p.~223, les deux conditions suivantes :

$\mathrm{F}1$) \deleted{\unsure{pour}} $\forall$ système \deleted{d'\unsure{objets}} de flèches \inserted{dans $M$}

\[
X(j)\xrightarrow{p(j)}\unsure{\varphi|\phi}(j)
\qquad\left|\begin{aligned}
&j\in\deleted{\unsure{I}}J_0=\mathrm{Ob}\,J\\
&\text{les }X(j)\text{ dans }\mathrm{Ob}\,C
\end{aligned}\right.
\]

il existe \inserted{un foncteur} $F:J\longrightarrow C$ et $F\xrightarrow{q}\varphi$, enfin un système de flèches $X(j)\xrightarrow{u(j)}F(j)$, \deleted{\unsure{rendant}} dans \inserted{les} $C/\varphi(j)$ i.e. rendant comm les diagrammes \unsure{i.e.,|i.e.}

\diagram{\begin{tikzcd}[ampersand replacement=\&,column sep=small]
X(j)\arrow[rr,"u(j)"]\arrow[dr,"p(j)"']\&\&F(j)\arrow[dl,"q(j)"]\\
\&\varphi(j)\&
\end{tikzcd}}

$\mathrm{F}3_{\unsure{\pi|\kappa}})$ Pour \deleted{tout} système $(F_\alpha)_{\alpha\in A}$ \deleted{\unsure{de foncteurs} $\in\underline{\mathrm{Hom}}(J,$} \inserted{$(\operatorname{card} A\leq\pi)$} d'éléments dans \deleted{les} $\underline{\mathrm{Hom}}(C,J)$, et des flèches $F_\alpha\xrightarrow{p_\alpha}\varphi$, dans $\underline{\mathrm{Hom}}(J,M)$, et pour $\deleted{\unsure{F}}G$ dans $\underline{\mathrm{Hom}}(J,\unsure{C|I})$, et $\deleted{\unsure{F}}G\xrightarrow{\deleted{\unsure{p}}q}\varphi$ dans $\underline{\mathrm{Hom}}(J,M)$, \unsure{étant donné|dont donné} un système de flèches $\unsure{u_\alpha(j)}$.

% slanted
\margin{en fait, il suffit de vérifier $\mathrm{F}3$ (cf plus bas 4.9.7.1), i.e. \unsure{lors|cas|lorsq.} $A$ réduit à un élément}

\diagram{\begin{tikzcd}[ampersand replacement=\&,column sep=small]
F_\alpha(\unsure{j})\arrow[rr,"u_\alpha(j)"]\arrow[dr,"p_\alpha(j)"']\&\&G\unsure{_\alpha|{}}(\unsure{j})\arrow[dl,"q(j)"]\\
\&\varphi(j)\arrow[ul,phantom,"\text{comm.}"{near start,swap}]\&
\end{tikzcd}}

\margin{\deleted{\unsure{système}} \deleted{\unsure{flè}} $u_\alpha(j)$ flèche dans $C/\varphi(j)$}

il existe \deleted{$\unsure{H\xrightarrow{r}\varphi}$ \unsure{dans} $\underline{\mathrm{Hom}}(J,C)_{\unsure{/\varphi}}$} \inserted{$H\xrightarrow{\unsure{q|r}}\varphi\unsure{^{\prime}|{}}$ dans $\underline{\mathrm{Hom}}(J,\unsure{C})_{/\varphi}$} et une \deleted{\unsure{un}} flèche $G\xrightarrow{v}H$ dans $\underline{\mathrm{Hom}}(J,C)_{/\varphi}$, telle que pour \unsure{tt|tout} $\alpha\inserted{{}\in A}$, le système de flèches composées

\scanpage{311}
\mspage{22\deleted{\unsure{6}}7}
% archive sheet number: 310

\diagram{\begin{tikzcd}[ampersand replacement=\&,column sep=large]
F_\alpha(j)\arrow[r,"u_\alpha(j)"]\arrow[dr,dashed,"p_\alpha(j)"']\&G(j)\arrow[r,no head,"v(j)"]\arrow[d,dashed,"\deleted{\unsure{p}}q(j)"]\&H(j)\arrow[dl,dashed,"q(j)"]\\
\&\unsure{\varphi|\phi}(j)\arrow[ul,phantom,"\text{comm.}"{near start}]\arrow[ur,phantom,"\text{comm.}"{near start}]\&
\end{tikzcd}}

(pour $j\in\inserted{J_0=}\mathrm{Ob}\,J$) provienne d'un \uline{foncteur}

\[w_\alpha:F_\alpha\longrightarrow H.\]

\deleted{C'est donc \unsure{en}} Pour prouver 4.9.3, nous sommes donc ramenés à prouver les deux \unsure{propositions|propriétés} $\mathrm{F}1$), $\mathrm{F}3_{\unsure{\pi|\kappa}})$ \unsure{précédentes}, sous l'une des deux conditions a) ou b) de 4.9.3.

\deleted{4.9.6.5. \uline{Démonstration de} 4.9.6.3}

\uline{Corollaire} 4.9.6.5 (de 4.9.6.3). Soit $(f_\alpha:I_\alpha\unsure{\mathbin{-}|\longrightarrow} I'_\alpha)_{\alpha\in\unsure{A}}$ une famille de foncteurs cofinaux, les \deleted{\unsure{catégories}}$I_\alpha$ sont \deleted{\unsure{étant}} filtrantes. Alors le foncteur produit $f=\prod f_\alpha:I=\prod_\alpha I_\alpha\longrightarrow I'=\prod_\alpha I'_\alpha$ est cofinal.

En effet, \deleted{\unsure{On sait que $I$ est filtrante|On sait que ces foncteurs sont},} \inserted{il faut vérifier que les} \deleted{catégories} $ {}_{i'}\backslash I$, pour $i'\in\mathrm{Ob}\,I'$, sont $0$-connexes. Or si $i'=(i'_\alpha)_{\alpha\in A}$, on a

\[ {}_{i'}\backslash I\simeq\prod_\alpha({}_{i'_\alpha}\backslash I_\alpha),\]

et par 4.9.6.3, 1\textsuperscript{o}) $\unsure{a})\Rightarrow c)$, les $ {}_{i'_\alpha}\backslash I_\alpha$ sont filtrantes, donc leur produit est filtrant

\scanpage{312}
\mspage{22\deleted{\unsure{7}}8}
% archive sheet number: 311

(cf. 4.9.6.2), donc $0$-connexe, qed.

\uline{Corollaire} 4.9.6.6. Soi\deleted{\unsure{en}}t $\unsure{\pi|\kappa}$ un cardinal infini. Tout produit de catégories \unsure{grdes|gdes} devant $\pi$ est \unsure{grde|gde} devant $\pi$.

Soit $I=\prod_{\alpha\in A}I_\alpha$, soit $B$ un ens de cardinal $\leq\pi$, il faut prouver que $\deleted{\unsure{B}}I\overset{\inserted{\unsure{\mathrm{diag}_B I|\mathrm{diag}_B}}}{\unsure{\mathbin{-}|\longrightarrow}} I^B$ est cofinal. Or ce foncteur s'identifie au produit des foncteurs similaires $I_\alpha\mathbin{-} I_\alpha^B$, lesquels sont cofinaux, les catégories sources $I_\alpha$ étant de plus filtrantes. Donc on peut conclure par le cor. 4.9.6.5.

4.9.6.7. \uline{Démonstration de} 4.9.6.3,. \inserted{(Partie 1\textsuperscript{o})}

$a\Rightarrow b$ \quad On sait déjà que a) implique que $I'$ est filtrante (SGA 4 I 8.1.3. b). Reste à prouver que pour \unsure{tt|tout} $i'\in\mathrm{Ob}\,I'$, $ {}_{i'}\backslash I$ est filtrante.\textsuperscript{*} Deux \unsure{conditions|condition} $\mathrm{PF}1'$ et $\mathrm{PF}2$\unsure{'|,} à vérifier (en plus \unsure{des|du} fait que \deleted{les} $ {}_{i'}\backslash I\ne\varnothing$, qui est connu par cofinalité de $f$, \unsure{impliquant} que $ {}_{i'}\backslash I$ est $0$-connexe) :

% margin note written sideways
\margin{Court-circuit : on sait déjà que $ {}_{i'}\backslash I$ est $0$-connexe \deleted{\unsure{d'o}} d'autre part elle est \uline{loc.} \unsure{filtr.te|filtrante}, car coétale sur $I$ qui est \unsure{filtr.te}, donc elle est filtrante..}

1) Soient $(i,u')$, $(j,v')$ deux éléments de $ {}_{i'}\backslash I$, il faut voir qu'ils sont majorés par un troisième élément. Or $I$ étant filtrante, $i,j$ sont majorés par $k$. On aurait fini si on savait que $f(w_1)u'=f(w_2)v'$, \unsure{soit|Soit} $w'$, car alors $(k,w')$ ferait l'affaire.

\margin{\diagram{\begin{tikzcd}[ampersand replacement=\&,row sep=tiny,column sep=tiny]
\&i\arrow[rr,"w_1"]\arrow[dd,no head]\&\&k\arrow[dd,no head]\\
\&\&j\arrow[ur,"w_2"']\arrow[dd,no head,"\unsure{\sim|?}"']\&\\
i'\arrow[r,"u'"]\arrow[rrd,"v'"']\& f(i)\arrow[rr,dashed]\&\&f(k)\\
\&\&f(j)\arrow[ur]\&
\end{tikzcd}}}

\scanpage{313}
\mspage{\deleted{\unsure{2}}229}
% archive sheet number: 312

\diagram{\begin{tikzcd}[ampersand replacement=\&,row sep=small,column sep=small]
\&k\arrow[r,"w"]\arrow[d,no head]\&\overset{2}{\overline{k}^{\unsure{*|2}}}\arrow[d,no head]\\
i'\arrow[r,shift left]\arrow[r,shift right,no head]\&f(k)\arrow[r]\&f(\unsure{k'|\overline{k}})
\end{tikzcd}}

Mais, \unsure{en appliquant|on applique} $\mathrm{F}2$ sur $f$ à la double flèche en question \diagram{\begin{tikzcd}[ampersand replacement=\&,cramped,column sep=small] i'\arrow[r,shift left=.6ex,no head]\arrow[r,shift right=.6ex,no head]\&f(k) \end{tikzcd}}, on trouve $w:k\longrightarrow\overline{k}$ telle que

\[
\underset{\substack{\Vert\\ f(wu_1)u'}}{f(w)(f(u_1)u')}=
\underset{f(ww_1)v'}{\underbrace{f(w)(f(\unsure{w_2|w_1})v')}}\;,
\]

et remplaçant $k$ par $\overline{k}^{\unsure{*}}$, et $u_1,v_1$ par $wu_1,wv_1$, on trouve ce qu'on veut.

2) PF(2) pour $ {}_{i'}\backslash I$, i.e. on se donne un diagramme ci-contre, avec $f(u_1)u'=f(u_2)u'$, \deleted{\unsure{\(\xi\)}} \inserted{soit $v'$,} (double-flèche dans $ {}_{i'}\backslash I$ de $(i,u')$ dans $(j,v')$), \deleted{\unsure{On}} veut l'égaliser à droite. Or

\diagram{\begin{tikzcd}[ampersand replacement=\&,column sep=tiny,row sep=small]
\&i\arrow[r,shift left,"u_1"]\arrow[r,shift right,"u_2"']\arrow[d,no head]\&j\arrow[r,dashed,"w"]\arrow[d,no head]\&j'\\
i'\arrow[r,"u'"]\&f(i)\arrow[r,shift left,"f(u_1)"]\arrow[r,shift right,"f(u_2)"']\&f(j)\arrow[r,dashed,no head,"f(w')"]\&f(j')
\end{tikzcd}}

$I$ étant \unsure{filtrant|filtrante}, il existe $w:j\unsure{\mathbin{-}|\longrightarrow}j'$ telle que $wu_1=wu_2$, or cette flèche définit \deleted{\unsure{s}} une flèche $(j,v')\longrightarrow(j',f(w')v')$, qui égalise la double flèche donnée.

$b\Rightarrow b'$ est clair

$\unsure{b')|b)}\Rightarrow a)$. Comme les $ {}_{i'}\backslash I$ sont filtrantes, donc $0$-connexes, $f$ est cofinal, il reste à prouver que $I$ est filtrante, \deleted{\unsure{donc}} à vérifier \deleted{\unsure{qui est}} $\mathrm{PF}0$ ($I$ \deleted{\unsure{0}}connexe) $\mathrm{PF}1$, $\mathrm{PF}2$ pour $I$. $I$ est $0$-connexe, car $I'$ l'est \deleted{(étant filtrante)}, et $f$ cofinal i.e. $W_0$-coasphérique, d'où $\mathrm{PF}0$.

\scanpage{314}
\mspage{230}
% archive sheet number: 313

\diagram{\begin{tikzcd}[ampersand replacement=\&,row sep=tiny,column sep=small]
\&i_1\arrow[dr,dotted,no head,"v_1"]\&\\
i_0\arrow[ur,"u_1"]\arrow[dr,"u_2"']\&\&i_3\\
\&i_2\arrow[ur,dotted,"v_2"']\&
\end{tikzcd}}

Prouvons $\mathrm{PF}\deleted{\unsure{2}}1$, i.e. que, si

\diagram{\begin{tikzcd}[ampersand replacement=\&,row sep=tiny,column sep=small]
\&i_1\\ i_0\arrow[ur,no head]\arrow[dr,no head]\&\\ \&i_2
\end{tikzcd}}

est donné, on peut compléter en un carré commutatif.

\diagram{\begin{tikzcd}[ampersand replacement=\&,row sep=tiny,column sep=small]
\&f(i_1)\arrow[dr,"f(\unsure{v_1|v'_1})"]\&\\
f(i_0)\arrow[ur,no head,"f(u_1)"]\arrow[dr,no head,"f(u_2)"']\&\&f(i'_3)\\
\&f(i_2)\arrow[ur,no head,"f(v'_2)"']\&
\end{tikzcd}}

\deleted{Mais $f(u_1),f(u_2)$ se complètent dans $I'$ en un carré commutatif sur un $i'\in\mathrm{Ob}\,I'$, et comme $f$ est cofinal, \unsure{OPS|OPG} $i'$ de la forme $f(i_3)$ ; mais les flèches $v'_1,v'_2$ ne proviennent peut-être pas de flèches $v_1,v_2$.}

Mais considérons les objets $\unsure{\zeta|\xi}_1=(\unsure{i_1|i_0},f(u_1)),\zeta_2=(i_2,f(u_2))$ de $ {}_{f(i_0)}\backslash I$, et utilisant $\mathrm{PF}1'$ pour cette catégorie, on trouve $\zeta_3=(i_3,w')$ dans ladite catégorie, et des flèches de $\zeta_1,\zeta_2$ en $\zeta_3$ i.e. $v_1:i_1\unsure{\mathbin{-}|\longrightarrow}i_3$, telles que

\[
\begin{aligned}
\underbrace{f(v_1)f(u_1)}_{f(v_1u_1)}&=w'\deleted{\unsure{x_2}},\quad
\underbrace{f(v_2)f(u_2)}_{f(\unsure{w_2|v_2}u_2)}=w',\qquad\text{donc}\\
f(\unsure{e_1|v_1}u_1)&=f(v_2u_2).
\end{aligned}
\]

Ainsi, on a une double flèche \diagram{\begin{tikzcd}[ampersand replacement=\&,cramped,column sep=small] i_0\arrow[r,shift left=.65ex,"v_1u_1"]\arrow[r,shift right=.65ex,"v_2u_2"']\&i_3 \end{tikzcd}}, égalisée par $f$. Si on prouve $\mathrm{PF}2$ sur $I$, on trouve $i_3\xrightarrow{\alpha}\overline{i}_3^{\unsure{*}}$ telle que $\alpha(v_1u_1)=\alpha(v_2u_2)$ i.e. $(\alpha v_1)u_1=(\alpha v_2)u_2$, et quitte à remplacer $\unsure{i'_3|i_3}$ par $\overline{i}_3$, $v_1,v_2$ par $\alpha v_1,\alpha v_2$, on a \unsure{gagné}

Reste \deleted{à} : prouver $\mathrm{PF}2$ sur $I$, i.e à égaliser une double-flèche

\diagram{\begin{tikzcd}[ampersand replacement=\&]
i_0\arrow[r,shift left,"u_1"]\arrow[r,shift right,"u_2"']\&i_1,
\end{tikzcd}}

\scanpage{315}
\mspage{231}
% archive sheet number: 314

\diagram{\begin{tikzcd}[ampersand replacement=\&,row sep=small,column sep=small]
i_0\arrow[r,shift left,"u_1"]\arrow[r,shift right,"u_2"']\arrow[d,no head]\&i_1\arrow[r,dotted,"\deleted{\unsure{w}}"]\&i_2\\
f(i_0)\arrow[r,shift left,"f(u_1)"]\arrow[r,shift right,"f(u_2)"']\& f(i_1)\arrow[r,dotted,"v'"]\& f(i_2)
\end{tikzcd}}

\deleted{On \unsure{considère les objets}} \inserted{$I'$ satisfaisant à $\mathrm{PF}2$, on}

\deleted{$\unsure{\zeta|\xi}_0=(i_0,\mathrm{id}_{f(i_0)}),\ \unsure{\zeta}$ $\zeta_1=(i_1,f(u_1)),\zeta_2=(i_1,\unsure{f(u_2)}))$ de $ {}_{f(i_0)}\backslash I$, comme cette cat. est filtrante, on trouve un objet \deleted{\unsure{qui}} $\zeta_3=(i_2,v')$ \unsure{qui|que} les \unsure{majore|majorant}, donc \deleted{\unsure{deux}} flèches \deleted{\unsure{de}} \diagram{\begin{tikzcd}[ampersand replacement=\&,cramped,column sep=small] i_1\arrow[r,shift left=.65ex,"v_1"]\arrow[r,shift right=.65ex,"v_2"']\&i_2 \end{tikzcd}} telles que $f(v_1)f(u_1)=f(v_2)f(u_2)$ i.e. $f(v_1u_1)=f(v_2u_2)$.}

peut égaliser la double flèche $f(u_1),f(u_2)$ par une flèche $f(i_1)\overset{v'}{\unsure{\mathbin{-}|\longrightarrow}}i'$ dans $I'$, et comme \deleted{les} $ {}_{i'}\backslash I$ non vide, OPS que $i'$ est de la forme $f(i_2)$. \deleted{Soit}

\deleted{$w'\overset{\unsure{\text{déf}}}{=}v'f(u_1)=v'f(u_2)$,}

\deleted{de \unsure{considérer}} Admettons pour un instant la condition $\mathrm{F}3$ de c), que nous allons appliquer à $(i_1,i_2,\deleted{\unsure{u}}v')$. Il existe donc $i_2\xrightarrow{v}\overline{i}_2$, $w:i_{\unsure{1}}\mathbin{-}\overline{i}_2$ telles que l'on ait $f(\unsure{w|v})v'=f(w)$. Il s'ensuit que $f(w)$ égalise $f(u_1)$ et $(fu_2)$, i.e. qu'on a $f(w\unsure{u_1})=f(wu_2)$, i.e. la double-flèche initiale est \unsure{remplacée|renforcée} par $(\overline{u}_1=wu_1,\overline{u}_2=wu_2)$, telle que $f(\overline{u}_1)=f(\overline{u}_2)$. En somme, on est ramené au cas \deleted{\unsure{d'}}une double-flèche \deleted{\unsure{flèches}} $(u_1,u_2)$ telle que $f(u_1)=f(u_2)$, \inserted{soit $v'$.} Mais alors soit $\zeta_1=(i_{\unsure{0|1}},v')$, \inserted{$\zeta_0=(i_0,\mathrm{id}_{f(\unsure{i_0|i'_0})})$,} \deleted{\unsure{la deux}} une double flèche $(u_1,\unsure{u_2|v_2}):$ \diagram{\begin{tikzcd}[ampersand replacement=\&,cramped,column sep=small] \zeta_0\arrow[r,shift left=.6ex]\arrow[r,shift right=.6ex]\&\zeta_1 \end{tikzcd}} dans $ {}_{f(i_0)}\backslash I$. Cette catégorie

\scanpage{316}
\mspage{232}
% archive sheet number: 315

est filtrante, une \deleted{cette} double-flèche est \unsure{égalisable} \inserted{à dr.}, ce qui implique aussitôt que la double flèche $(\unsure{v_1,v_2|u_1,u_2|v_1,u_2})$ \uline{dans $I$} est aussi égalisable à droite.

Donc il reste à prouver que b') implique la condition $\mathrm{F}3$. Mais considérons les deux objets $\unsure{\zeta|\xi}_0=(i_0,\mathrm{id}_{f(i_0)})$ $\zeta_1=(i_{\unsure{1}},u')$ de $ {}_{f(i_0)}\backslash I$,

\diagram{\begin{tikzcd}[ampersand replacement=\&,column sep=small]
i_0\arrow[rr,dashed,bend left=16,"w"]\arrow[d,no head]\&i_1\arrow[r,dashed,no head,"v"']\arrow[d,no head]\&i_2\arrow[d,no head]\\
f(i_0)\arrow[r,"u'"]\arrow[rr,dashed,bend right=18,no head,"v'"']\&f(\unsure{i_1})\&f(i_2)
\end{tikzcd}}

et utilisons la condition de filtration $\mathrm{PF}1$ sur cette catégorie, donc on trouve un objet $\zeta_{\unsure{2|1}}=(i_2,v')$, et deux flèches \deleted{$\unsure{\zeta_0\rightrightarrows\zeta_1}$}

\diagram{\begin{tikzcd}[ampersand replacement=\&,row sep=tiny]
\zeta_0\arrow[dr,no head,"w"]\&\\
\&\zeta_2,\\
\zeta_1\arrow[ur,no head,"v"']\&
\end{tikzcd}}

i.e. $i_0\overset{w}{\unsure{\mathbin{-}|\longrightarrow}}i_2$ et $i_1\overset{v}{\mathbin{-}}i_2$ telles qu'on ait $f(w)\mathrm{id}_{\unsure{j_0|i_0}}=f(v)u'=v'$, d'où $f(w)=f(\unsure{w|v})u'$, ce qui est la condition $\mathrm{F}3$.

\deleted{Prouvons} Ainsi on a établi

\[a\Leftrightarrow b\Leftrightarrow b',\]

prouvons maintenant

\uline{$a\Rightarrow c$.} On sait déjà que a) implique $I'$ filtrante, et il reste à établir $\mathrm{F}3$. \deleted{\unsure{Avec}} Or $I$ étant filtrante, il existe \deleted{\unsure{$i_2\longrightarrow$}} des flèches

\diagram{\begin{tikzcd}[ampersand replacement=\&,row sep=tiny]
i_0\arrow[dr,"w"]\&\\
\&i_2.\\
i_1\arrow[ur,"v"']\&
\end{tikzcd}}

\unsure{On|Or}

\scanpage{317}
\mspage{233}
% archive sheet number: 316

\unsure{n'a} pas forcément $f(w)=f(v)u'$, on a là une double flèche \deleted{\unsure{dans I}} $f(i_0)\unsure{=|\rightrightarrows}f(i_2)$ dans $I'$. Mais

\deleted{$\unsure{\zeta|\xi}$ \inserted{$ {}_{f(i_0)}\backslash I$} étant filtrante, cette double flèche peut être égalisée, les deux objets $\zeta'_2=(i_2,\deleted{\unsure{\zeta}}f(w)),\zeta''_2=(i_2,$}

par la condition $\mathrm{F}2$ sur $f$, il existe $i_2\xrightarrow{\alpha}\overline{i}_2$ dans $I$ telle que $f(\alpha)$ égalise ladite double flèche, donc $f(\underbrace{\alpha w}_{\overline{w}})=f(\underbrace{\alpha v}_{\overline{v}})u'$, et quitte à remplacer $i_2$ par $\overline{i}_2$, et \deleted{\unsure{s}} $w,v$ par $\overline{w}=\alpha w$, $\overline{v}=\alpha\unsure{w|v}$, on \unsure{gagne}.

Reste à prouver (pour la partie 1\textsuperscript{o}), que si $f$ est \uline{fidèle}, alors $c)\Rightarrow\deleted{\unsure{a}}b)$. Il reste donc à prouver que \unsure{b)|c)} implique que \deleted{\unsure{pour}} pour $i'$ dans $I'$, $ {}_{i'}\backslash I$ est filtrante. On sait déjà par $\mathrm{F}1$ que $ {}_{i'}\backslash I$ est non vide, il reste à prouver les conditions $\mathrm{PF}1'$, $\mathrm{PF}2$ pour $ {}_{i'}\backslash I$. Comme $I'$ \unsure{filtrant|filtrante}, déjà on peut compléter

\diagram{\begin{tikzcd}[ampersand replacement=\&,column sep=small,row sep=tiny]
\&f(i_1)\arrow[dr,dashed,no head]\&\\
i'\arrow[ur,no head]\arrow[dr,no head]\&\&\unsure{j''|j'}\\
\&f(i_2)\arrow[ur,dashed,no head]\&
\end{tikzcd}}

\diagram{\begin{tikzcd}[ampersand replacement=\&,row sep=tiny,column sep=tiny]
\&i_1\arrow[rr,dashed,no head,"v_1?"]\arrow[dd,no head]\&\&i_3\arrow[r]\arrow[dd,no head]\&\overline{i}_3\arrow[dd,no head]\\
\&i_2\arrow[urr,no head,"v_2?"']\arrow[dd,no head]\&\&\&\\
\&f(i_1)\arrow[rr,no head,"v'_1"]\&\&f(i_3)\arrow[r,no head]\&f(\unsure{i'_3|\overline{i}_3})\\
i'\arrow[ur,no head,"u'_1"]\arrow[r,no head,"u'_2"']\&f(i_2)\arrow[urr,"v'_2"']\&\&\&
\end{tikzcd}}

, et \deleted{\unsure{comme}} par $\mathrm{F}1$ OPS $j'$ de la forme $f(i_3)$. Appliquant $\mathrm{F}3$ à $i_1,i_3,v'$, OPS (quitte à remplacer $i_3$ par $\overline{i}_3$)

\scanpage{318}
\mspage{234}
% archive sheet number: 317

que $v'_1$ provient d'un $v_1:i_1\unsure{\mathbin{-}|\longrightarrow}i_3$. Appliquant à nouveau $\mathrm{F}3$ à $i_2,i_3,v'_2$, OPS que $v'_2$ également provient d'un $v_2:i_2\mathbin{-}i_3$. Alors $\unsure{\zeta|\xi}_3=(i_3,v'_1u'_1=v'_2u'_2)$ est un majorant commun de $\zeta_1=(i_1,u'_1),\zeta_{\unsure{2}}=(i_2,u'_2)$.

Prouvons enfin $\mathrm{PF}2$ pour $ {}_{i'}\backslash I$, \unsure{donc|dans} on a un diagramme ci-contre avec $f(u_1)u'=f(u_2)u'$, (double flèche dans $ {}_{i'}\backslash I$),

\diagram{\begin{tikzcd}[ampersand replacement=\&,column sep=tiny,row sep=large]
\&\deleted{\unsure{i}}i_0\unsure{)|\prime}\arrow[r,shift left,"u_1"]\arrow[r,shift right,"u_2"']\&i_1\arrow[rr,dotted,bend left=15,"w"]\&i_2\arrow[r,dotted]\&\overline{i}_2\\
i'\arrow[r,"u'"]\&f(i_0)\arrow[r,shift left,"f(u_1)"]\arrow[r,shift right,"f(u_2)"']\&f(i_{\unsure{1}})\arrow[r,dashed,"v'"]\&f(i_2)\arrow[r]\&f(\overline{i}_2)
\end{tikzcd}}

on veut égaliser dans $ {}_{i'}\backslash I$. Mais $I'$ étant \unsure{filtrant|filtrante}, on peut égaliser $(\deleted{\unsure{f}}f(u_1),f(u_2))$ par $f(i_{\unsure{1|3}})\longrightarrow j'$, et à cause de $\mathrm{F}1$ OPS $j'=f(i_2)$. Appliquant encore $\mathrm{F}3$ à $(i_1,i_2,v')$, on trouve $i_1\longrightarrow\overline{i}_2$ tel que $f(w)$ égalise $f(u_1),\unsure{fu_2)|f(u_2)}$ :

\[f(wu_1)=f(wu_2).\]

Mais $f$ étant \uline{fidèle}, cela implique $wu_1=wu_2$, et on a donc égalisé \unsure{la|les|des} \unsure{double-flèche|double-flèches} de $ {}_{i'}\backslash I$ par $w:(i_1,-)\longrightarrow(\overline{i}_2,-)$.

\uline{4.9.6.\unsure{8}.$^{\unsure{1}}$} \uline{NB} \inserted{\textcircled{$\uparrow$}} Si on ne suppose pas que $f$ soit fidèle, on voit que \deleted{\unsure{f}} c) implique \inserted{\unsure{pourtant|cependant}} que les $ {}_{i'}\backslash I$ \deleted{\unsure{satis}}\unsure{satisfont|sfnt.} $\mathrm{PF}1'$ (elles sont donc \deleted{0-connexes et} pseudo-filtrantes et $0$-connexes) \deleted{\unsure{donc}} que $f$ est cofinal. \deleted{On voit}

\scanpage{319}
\mspage{235}
% archive sheet number: 318

On voit, par les $=$ arguments, que $I$ \deleted{\unsure{est}}satisfait $\mathrm{PF}1'$ i.e. est ps. \unsure{filtr.te|filtrante}. Ainsi, si on introduit la condition

\tikz[baseline=(b.base)]\node[draw,ellipse,inner sep=2pt](b){$b''$}; $I'$ est filtrante, \deleted{\unsure{$I$ et}} les $ {}_{i'}\backslash I$ $(i'\in\mathrm{Ob}\,I')$ sont pseudofiltrantes \inserted{(i.e. satisfont $\mathrm{PF}1'$)} (\deleted{donc} $f$ cofinal).

On trouve les implications

\[(a\Leftrightarrow b\Leftrightarrow b')\Rightarrow(b''\Leftrightarrow c)\]

et \deleted{\unsure{que}} les conditions équivalentes c), b'') impliquent que $I$ est \inserted{également} ps. filtrante, et que pour \unsure{tt|toute|$\forall$} double-flèche \inserted{\diagram{\begin{tikzcd}[ampersand replacement=\&,cramped,column sep=small] \unsure{\zeta|\xi}_0\arrow[r,shift left=.6ex,"u_1"]\arrow[r,shift right=.6ex,"u_2"']\&\zeta_1 \end{tikzcd}}} dans un $ {}_{i'}\backslash I$ ou dans $I$, il existe une flèche $\zeta_1\overset{v}{\longrightarrow}\zeta_2$ dans la $=$ catégorie, telle que $vu_1$ et $vu_2:$ \diagram{\begin{tikzcd}[ampersand replacement=\&,cramped,column sep=small] \zeta_0\arrow[r,shift left=.6ex]\arrow[r,shift right=.6ex]\&\zeta_2 \end{tikzcd}} aient \uline{$=$ image} \uline{dans $I'$}.

Donc si $f$ est fidèle, alors ces conditions \unsure{a)|c)} c' impliquent que $I$ et les $ {}_{i'}\backslash I$ satisfont aussi la condition $\mathrm{PF}2$ de la double flèche, donc sont filtrantes, \deleted{donc} dans ce cas les cinq conditions a) b) b'), c) c') sont équivalentes. Mais aussi lorsque $I$ est ordonnée, \deleted{\unsure{donc|de}} \unsure{sorte|suite|sûr} que la condition de la double flèche dans $I$ est satisfaite automatiquement.

\scanpage{320}
\mspage{236}
% archive sheet number: 319

2\textsuperscript{o}) Notons aussi que dans la condition b'), on ne peut supprimer l'hypothèse que $I'$ satisfait $\mathrm{PF}2$ (en plus de $I'$ $0$-connexe et les $ {}_{i'}\backslash I$ filtrantes). Pour le voir, prenons \deleted{$f$} $I'=I$, $f=\mathrm{id}$, donc l'hyp. \unsure{b')|a)} $(\Leftrightarrow b\Leftrightarrow b')$ signifie simplement que $I$ est filtrante, $f$ étant bien sûr cofinal.

Mais b') affaiblie signifie que $I$ est $0$-connexe et les $ {}_{i}\backslash I$ sont filtrantes, et on sait que ceci n'implique \uline{pas} que $I$ soit filtrante, $=$ si on suppose $I$ pseudo-filtrante. Un \deleted{exemple} simple est donné par $I$ définie par

\[
\mathrm{Ob}\,I=\mathbb{N},\qquad
\deleted{\unsure{\mathrm{Fl}}}\deleted{\unsure{\mathrm{Ens}}}
\mathrm{Hom}_{I}(i,j)=
\begin{cases}
\varnothing&\text{si }i>j\\
\{\mathrm{id}_i\}&\text{si }i=j\\
\{\pm1\}_{\unsure{i,j|j,i}}&\text{si }i<j,
\end{cases}
\]

la composition des flèches se faisant par la multiplication des signes $\unsure{\varepsilon|\epsilon}_{k,j}\varepsilon'_{ji}=(\varepsilon\unsure{\bar{\varepsilon}'|\varepsilon'})_{ki}$. (en posant $\mathrm{id}_i=(+1)_{i,i}$). Les catégories $ {}_i\backslash I$ sont \inserted{\unsure{isom.} entre elles et sont} \uline{ordonnées} et ps filtr. donc filtrantes, \deleted{\unsure{d'ailleurs}} puis $I$ est ps. filtrante, \unsure{pour}\deleted{\unsure{tant}} elle n'est pas filtrante, car ne satisfait pas $\mathrm{PF}\unsure{2|?}$.

3) D'autre part, si $f$ n'est pas fidèle, alors c) (équivalente à \uline{b'')}) n'implique pas a), comme on voit en prenant $I'=e$ ; alors b'') équivaut à $I$ ps. filtrante, qui n'implique pas que $I$ soit filtrante.

\scanpage{321}
\mspage{237}
% archive sheet number: 320

\uline{4.9.6.9}. Démonstration de 4.9.6.3, \uline{partie 2\textsuperscript{o}}.

C'est essentiellement ``la $=$'' démonstration que pour la partie 1\textsuperscript{o}).

$a_{\unsure{\pi|\kappa}})\Rightarrow b_\pi)$. On peut court-circuiter la démonstration directe, en regardant, pour un ensemble $A$ de cardinal $\leq\pi$, le diagramme

\diagram{\begin{tikzcd}[ampersand replacement=\&]
I\arrow[r,"\delta_A"]\arrow[d,"f"']\&I^A\arrow[d,"f^A"]\\
I'\arrow[r,"\delta_{A'}"']\&I'^{A}.
\end{tikzcd}}

Comme $f$ est cofinal, \deleted{\unsure{$f$}} $I$ filtrante, $f^A$ est aussi cofinal (4.9.\unsure{8|6}.5), d'autre part $I$ \unsure{grde|gde} devant $\pi$ signifie que $\delta_A$ est cofinal ($\forall\, A$ avec $\operatorname{card} A\leq\pi$). Donc $f^A\delta_A=\delta_{A'}f$ est cofinal, donc $\delta_{A'}f$ et $f$ le sont, donc $\delta_{A'}$ aussi, ce qui signifie que $I'$ est grde devant $\pi$.

Reste à prouver que pour $i'$ dans $I'$, $ {}_{i'}\backslash I$ est grde devant $\pi$. Cela résulte du fait que $ {}_{i'}\backslash I\xrightarrow{\unsure{q|\varphi}}I$ est coétale, \inserted{que $ {}_{i'}\backslash I$ est $0$-connexe,} et du

\uline{Lemme} 4.9.6.9.1. Soit $I$ \unsure{une} catégorie filtrante (\inserted{resp.} grde devant $\pi$). \deleted{\unsure{et}} Alors pour \unsure{tte|tout} $J$ sur $I$ coétale sur $I$, \inserted{et $0$-connexe,} $J$ est filtrante (resp. grde devant $\pi$).

\deleted{Il faut vérifier} \deleted{\unsure{qui résulte}}

Le cas filtrant est déjà connu (cf II, \ldots). Pour le cas ``grd devant $\pi$'', on va utiliser le

\uline{\deleted{\unsure{Lemme}}Corollaire} 4.9.6.9.2. Soit $J$ une catégorie. Pour que $J$ soit grde devant $\pi$, il

\scanpage{322}
\mspage{238}
% archive sheet number: 321

f et s. qu'elle \inserted{soit $0$-connexe, et} satisfasse $\deleted{\unsure{P}}\mathrm{PF}_{\unsure{\pi|\kappa}}2$ (cf page 81) ainsi que la variante suivante de $\mathrm{PF}_\pi1'$

$\mathrm{PF}_\pi1)$ Pour \deleted{\unsure{de|un}}tout objet $\deleted{\unsure{y}}x_0$ de $J$, $ {}_{x_0}\backslash J$ satisfait $\mathrm{PF}_\pi1'$, i.e. pour \unsure{tte|toute} famille de flèches $x_0\xrightarrow{u_\alpha}\deleted{\unsure{x}}y_\alpha$ $(\alpha\in A,\operatorname{card} A\leq\pi)$, existe une \deleted{\unsure{des}} flèche $x_0\xrightarrow{v}\deleted{\unsure{y}}z$ et des flèches $v_\alpha:y_\alpha\unsure{\mathbin{-}|\longrightarrow}z$, avec $v_\alpha u_\alpha=v$ $\forall\,\alpha\unsure{\in|+} A$.

Démonstration. Si $J$ est \unsure{grde|gde} devant $\pi$, alors \deleted{\unsure{alors}} $\mathrm{PF}_\pi1)$ résulte aussitôt de la conjonction de $\mathrm{PF}_\pi1'$ et $\mathrm{PF}_\pi2$, d'où le ``il faut''. Inversement, supposons que $J$ satisfasse $\mathrm{PF}_\pi1,\mathrm{PF}_\pi2$, et soit $0$-connexe, il faut prouver qu'elle satisfait $\mathrm{PF}_\pi1'$. Mais on sait déjà que les conditions en question, \inserted{qui} impliquent \deleted{que} \inserted{$\mathrm{PF}1$, $\mathrm{PF}2$, $0$-conn.,} \deleted{que} \inserted{impliquent que} $J$ est filtrante. Donc si on a une famille $(x_\alpha)_{\alpha\in A}$ \deleted{non vide}, $\operatorname{card} A\leq\pi$, choisissant \inserted{un} $x_0\in\mathrm{Ob}\,J$, on peut pour \unsure{tt|tout} $\alpha$ trouver $y_\alpha$ qui majore $x_0$ et $x_\alpha$, d'où des flèches

\diagram{\begin{tikzcd}[ampersand replacement=\&,row sep=tiny]
x_\alpha\arrow[dr,dotted]\&\\
\&y_\alpha,\\
x_0\arrow[ur]\&
\end{tikzcd}}

on peut \deleted{\unsure{leur}} appliquer $\mathrm{PF}_\pi1$, d'où des flèches $y_\alpha\longrightarrow z$, et en composant des $x_\alpha\mathbin{-}z$, qed.

\deleted{Corollaire} Prouvons maintenant le lemme. Il \deleted{\unsure{en}} suffit \inserted{(par le cor. précédent)} de prouver que $J$ satisfait à $\mathrm{PF}_\pi1$ et $\mathrm{PF}_\pi2$. Or $\mathrm{PF}_\pi1$ est immédiat, puisque $\forall\,x_0$ de $J$, $ {}_{x_0}\backslash J\xrightarrow{\sim}{}_{\unsure{y_0|\underline{x}_0}}\backslash I$, où $y_0$ est

\scanpage{323}
\mspage{239}
% archive sheet number: 322

l'image de $x_0$ dans $I$. Quant à $\mathrm{PF}_{\unsure{\pi|\kappa}}2$, si on a des \diagram{\begin{tikzcd}[ampersand replacement=\&,cramped,column sep=small] x_0 \arrow[r,shift left=1.2ex,"u_\alpha"] \arrow[r] \arrow[r,shift right=1.2ex] \& x_1 \end{tikzcd}} dans $J$, d'où des \deleted{\unsure{$g(x_1)$|$g(x_0)$}} \diagram{\begin{tikzcd}[ampersand replacement=\&,cramped,column sep=small] g(x_0) \arrow[r,shift left=1.2ex,"g(u_{\unsure{1|\alpha|2}})"] \arrow[r,no head] \arrow[r,shift right=1.2ex] \& g(x_1) \end{tikzcd}}, on égalise \inserted{dans $I$} par $g(x_1)\xrightarrow{w}y$, mais comme $g$ est coétal, $\exists$ $x_1\xrightarrow{v}x_2$ dans $J$ telle que $g(v)=w$, donc on a $g(vu_\alpha)=g(\unsure{w|v}u_\beta)$ pour $\alpha,\beta\in A$, et comme $g$ est fidèle, on en conclut $vu_\alpha=vu_\beta$, \deleted{qed.} donc $v$ égalise $(u_\alpha)_{\alpha\in A}$, qed.

\deleted{\uline{Corollaire} 4.9.6.9.3. Soit $I$ une catégorie, $\pi$ un cardinal infini. Pour que $I$ soit \unsure{grde|gde} devant $\pi$, il f. et s. que \deleted{\unsure{$I$ soit $0$-connexe|$I$ soit connexe}} pour tout $i\in\mathop{\mathrm{Ob}}\nolimits I$, $ {}_i\backslash I$ le soit, \uline{et} que $I$ soit $0$-connexe.}

% slanted
\margin{Faux déjà dans le cas filtrant sans plus, cf \unsure{ex.|Ex.} page 236}

\deleted{C'est nécessaire par 4.9.6.9.1, puisque les $ {}_i\backslash I$ sont $0$-connexes et coétales sur $I$. C'est suffisant, car \deleted{\unsure{les|le cas}} cette condition implique que les $ {}_i\backslash I$ satisfont $\mathrm{PF}_\pi 1$, et par \deleted{\unsure{la|le}} 4.9.6.9.2 il \deleted{suffit} reste à vérifier \uline{$\mathrm{PF}_\pi 2$ pour $I$}.}

Prouvons $b_\pi)\longrightarrow a_\pi)$ : Si $I'$ \deleted{\unsure{est|est grde}} et les $ {}_{i'}\backslash I$ sont \unsure{grdes|gdes} devant $\pi$, alors $I$ l'est aussi (et $f$ est cofinal, les $ {}_{i'}\backslash I$ étant $0$-connexes, car filtrantes). Comme on sait déjà que $I$ est $0$-connexe ($I'$ l'étant et $f$ cofinal), il reste, par 4.9.6.9.2, à vérifier $\mathrm{PF}_\pi 1$ et $\mathrm{PF}_\pi 2$.

\scanpage{324}
\mspage{240}
% archive sheet number: 323

La démonstration donnée dans 1° pour le cas similaire \unsure{filtrant|filtrante} (\deleted{pour} \unsure{implication} $b'\Longrightarrow a)$, pages 231-232), s'applique mutatis mutandis. (La \deleted{condition} \inserted{propriété} $\mathrm{F}3$ \deleted{\unsure{pour}} a déjà été établie \deleted{dans} la partie 1°), comme conséquence de $b$, et peut donc être utilisée librement).

Donc on a prouvé $a_{\unsure{\pi|\kappa}}\Longleftrightarrow b_\pi$. Prouvons $a_\pi\Longrightarrow c_\pi$. On sait déjà que $a_\pi$ implique que $I'$ est \unsure{grde|gde} devant $\pi$, et bien sûr la condition $\mathrm{F}1$, reste à établir la validité de la condition $\mathrm{F\unsure{B|3}}_\pi$.

\diagram{\begin{tikzcd}[ampersand replacement=\&,column sep=small,row sep=small]
i_\alpha \arrow[rr,bend left=27,no head,"w_\alpha^0"] \arrow[d,no head] \& i \arrow[r,"v_\alpha^0"] \arrow[d,no head] \& j_\alpha \arrow[r,"\lambda_\alpha"] \arrow[d,no head] \& \unsure{g|j|h} \arrow[r,no head] \arrow[d,no head] \& \deleted{\unsure{h|k}} \arrow[d,no head] \\
f(i_\alpha) \arrow[r,"u'_\alpha"] \& f(i) \arrow[r,"f(v_\alpha)"] \& f(j_\alpha) \arrow[r,no head] \& f(\unsure{g|j|h}) \arrow[r,no head] \& f(\unsure{h|j'})
\end{tikzcd}}

Par $\mathrm{F}3$, $\forall\alpha\in A$ on peut trouver $j_\alpha$ dans $I$, et $w_\alpha^0:i_\alpha\mathbin{\unsure{-|\to}}j_\alpha$, $v_\alpha^0:i\to j_\alpha$ tels que

\[
f(v_\alpha)u'_\alpha=f(w_\alpha).
\]

Comme $I$ satisfait $\mathrm{PF}1^{\unsure{*|'}}_\pi$, \deleted{il existe} (cf 4.9.6.9.2) il existe \deleted{\unsure{un}} $h$ dans $I$ et des $\lambda_\alpha:j_\alpha\mathbin{-}\unsure{h|g}$, tels que $\lambda_\alpha v_\alpha^0=\lambda$ soit indépendant de $\alpha$. Alors \deleted{\unsure{les}} on a pour \unsure{tt|tous} $\alpha\in A$

\[
f(\lambda)u'_\alpha\left[=f(\lambda_\alpha)\underbrace{f(\unsure{w_\alpha|v_\alpha})u'_\alpha}_{f(w_\alpha)}\right]
=f(\underbrace{\lambda_\alpha w_\alpha^0}_{w_\alpha}),
\]

donc \deleted{\unsure{posons|juxtaposé}} $\lambda:i\mathbin{-}j$ et les \deleted{$v_\alpha=\lambda_\alpha v_\alpha^0:i\mathbin{-}j$} $w_\alpha=\lambda_\alpha w_\alpha^0:i_\alpha\longrightarrow j$ satisfont aux conditions voulues.

$c_\pi\deleted{\unsure{\Leftarrow}}\unsure{\Longrightarrow|\Longleftrightarrow}b_\pi''$ Il faut prouver que si $c_\pi$ est satisfait, alors les $ {}_{i'}\backslash I$ satisfont à $\mathrm{PF}1'_\pi$.

